\documentclass[11pt]{article}
\usepackage[T1]{fontenc}
\usepackage{lmodern}
\usepackage{amsmath,amssymb,amsthm}
\usepackage[margin=1in]{geometry}
\usepackage{microtype}
\usepackage{xurl}
\usepackage[hidelinks,breaklinks]{hyperref}
\hypersetup{pdftitle={Quantitative reconstruction from finite causal orders in a conformal diamond},pdfauthor={Daniel Eric Fredriksen},pdfsubject={Finite realizer rank rigidity and a restricted inverse stability bound},pdfkeywords={causal order, spacetime reconstruction, partial orders, Lean 4, inverse stability}}
\newtheorem{theorem}{Theorem}[section]
\newtheorem{lemma}[theorem]{Lemma}
\newtheorem{proposition}[theorem]{Proposition}
\newtheorem{corollary}[theorem]{Corollary}
\theoremstyle{remark}
\newtheorem{remark}[theorem]{Remark}
\newcommand{\R}{\mathbb R}
\newcommand{\K}{\mathcal K}
\newcommand{\TV}{\operatorname{TV}}
\newcommand{\dc}{d_{\mathrm{conf}}}
\newcommand{\rank}{\operatorname{rank}}
\newcommand{\norm}[1]{\left\lVert #1\right\rVert_\infty}
\title{Quantitative reconstruction from finite causal orders\\in a conformal diamond}
\author{Daniel Eric Fredriksen\\Quantyra Inc.\\\url{https://github.com/Quantyra/quantyra-null-cone-lean}}
\date{7 October 2026 \\ Unpublished literature revision of version 0.3.0}

\begin{document}
\maketitle
\begin{abstract}
Consider smooth volume densities on the unit square, bounded between
$1/2$ and $3/2$, with uniform marginals and Euclidean Lipschitz constant at
most two. In null coordinates these define a restricted class of
$1+1$-dimensional conformal diamonds with the same chronological relation.
Observing only the partial orders induced by independent volume samples,
we derive an explicit inverse estimate: the supremum difference of the
densities, modulo exchange of the null axes, is at most
$100(N^{-1/12}+\Delta_N)$, where $\Delta_N$ is the largest total variation
distance between the finite order laws through size $N$. A logarithmic
grid scale further gives $130((\log N/N)^{1/6}+\Delta_N)$ for every
$N\ge2$. The argument
uses an occupied-grid forcing construction that constrains every
two-order realizer, after one common axis exchange, to have interior
rank error at most $30rn$. Lean~4 proofs now cover this finite theorem,
the original inverse estimate, and all-size identifiability, including
the probability and analytic bridges and equality of labeled and
unlabeled order-law total variation. The logarithmic-grid improvement
and proper-time comparison remain ordinary mathematical proofs.
The constants are conservative, the exponents are not claimed optimal,
and originality remains provisional. Public source and reproducibility
artifacts support decentralized informal feedback and downstream testing.
\end{abstract}

\noindent\textbf{Keywords:} causal order; conformal geometry; two-dimensional
posets; quantitative identifiability; formal verification.\\
\textbf{Artifacts:} Lean source, pinned dependencies, verification audit,
and this manuscript are available in the public repository above.\\
\textbf{Revision status:} unpublished; no new DOI.\\
\textbf{Published baseline DOI:} version 0.3.0,
\href{https://doi.org/10.5281/zenodo.23214579}{doi:10.5281/zenodo.23214579}.\\
\textbf{Previous manuscript:} version 0.2.0,
\href{https://doi.org/10.5281/zenodo.23206773}{doi:10.5281/zenodo.23206773}.\\
\textbf{Software archive:} historical version 0.1.0,
\href{https://doi.org/10.5281/zenodo.23202763}{doi:10.5281/zenodo.23202763}.
The later full formalization is identified in Section~\ref{sec:formal}.\\
\textbf{License:} manuscript CC-BY-4.0; software Apache-2.0.

\section{Introduction and scope}
Chronology and volume have long motivated spacetime reconstruction.
Janson's equivalence theory for random posets supplies qualitative
identification of order kernels from all finite laws~\cite{Janson}.
Braun proves a smooth reconstruction theorem from chronological
adjacency laws under geometric hypotheses in spacetime dimension at
least three~\cite{Braun}. Henson gives a numerical construction of
approximate flat embeddings in two dimensions~\cite{Henson}.
These precedents establish the general reconstruction objective.

The present draft concerns an explicit finite inverse modulus for a
fixed, regular, two-dimensional class. Orientation forcing and modular
decomposition are established methods~\cite{KlavikZeman}. Winkler's
\emph{Random orders}, Theorem~7 and its proof, use existential poset
configurations to recover finite coordinate-order patterns up to axes
in the infinite box-dense setting~\cite{Winkler1985}. In the flat
two-dimensional random model, the typical-event argument in
\emph{Random orders of dimension 2}, Section~II and the proof of
Theorem~4.3, leaves only disjoint adjacent twin transpositions and one
global axis exchange~\cite{Winkler}. It follows that every realizer has
rank error at most one after that exchange, a stronger flat-model
conclusion than our $30rn$ estimate. Exact unique realizability is
unnecessary for this recovery.

Our candidate contribution is narrower: explicit deterministic
occupied-grid quantification and its uniform statistical use over the
nonconstant positive Lipschitz class, leading to a finite
law-to-coefficient inverse modulus. The independent uniform-permutation
counting in Winkler's finite model does not immediately give this
uniform conclusion for dependent coordinate rankings. We have not
located an equivalent or directly implying inverse bound in the
inspected literature. The full-text comparison and its remaining
structural/citation gaps are recorded in the repository; originality
remains provisional, with no certified priority claim.

The result concerns geometric inference within a stipulated class.
It supplies neither new physical dynamics nor reconstruction of an
arbitrary spacetime. No curvature estimate, practical sample-complexity
claim, or societal carrying-capacity consequence is established.

\section{Model and main estimate}
Let $D=[0,1]^2$, with null coordinates $(u,v)$, and set
\[
 g_\rho=-\rho(u,v)(du\otimes dv+dv\otimes du).
\]
The interior of $D$ is the spacetime; its closure is used for coefficient
norms and endpoint limits. Its volume measure is $\rho\,du\,dv$, and
chronology is
\[
 (u,v)\prec (u',v')\quad\Longleftrightarrow\quad u<u',\ v<v'.
\]
Let $\K$ consist of densities smooth on a neighborhood of $D$ such that
\begin{equation}\label{eq:class}
 \tfrac12\le\rho\le\tfrac32,\qquad
 |\rho(x)-\rho(y)|\le2|x-y|_2,\qquad
 \int_0^1\rho(u,v)\,dv=\int_0^1\rho(v,u)\,dv=1
 \quad(u\in[0,1]).
\end{equation}
In particular, total volume is one. Write $T(u,v)=(v,u)$ and
$\rho^T=\rho\circ T$. The uniform-marginal gauge is fixed; the residual
equivalence used here is axis exchange, not time reversal.

For $k$ independent samples from $\rho\,du\,dv$, let $\nu_k(\rho)$ be
the law of their labeled strict partial order. Labels are event indices,
not coordinate rankings. For $N\ge2$, define
\[
 \Delta_N(\rho,\sigma)=\max_{2\le k\le N}
       \TV(\nu_k(\rho),\nu_k(\sigma)),\qquad
 \dc(\rho,\sigma)=\min\{\norm{\rho-\sigma},\norm{\rho-\sigma^T}\}.
\]
Total variation is $\sup_A|\mu(A)-\nu(A)|$.

\begin{theorem}[Inverse estimate]\label{thm:main}
For every $\rho,\sigma\in\K$ and integer $N\ge2$,
\begin{equation}\label{eq:main}
 \dc(\rho,\sigma)\le100\bigl(N^{-1/12}+\Delta_N(\rho,\sigma)\bigr).
\end{equation}
\end{theorem}
The theorem is proved in Sections~\ref{sec:probability}--\ref{sec:inverse}
using the finite theorem of Section~\ref{sec:finite}. The complete
original estimate and Corollary~\ref{cor:identification} are now formalized
for both labeled and unlabeled directed-order laws; see
Section~\ref{sec:formal}.

\begin{corollary}[All-size identification]\label{cor:identification}
If $\nu_k(\rho)=\nu_k(\sigma)$ for every $k$, then
$\rho=\sigma$ or $\rho=\sigma^T$ on $D$.
\end{corollary}
\begin{proof}
Every $\Delta_N$ vanishes. Let $N\to\infty$ in~\eqref{eq:main}.
The minimum of the two fixed nonnegative supremum norms is zero.
\end{proof}

\section{Occupied grids and finite rank rigidity}\label{sec:finite}
Take $n\ge1$ points $x_i=(u_i,v_i)\in(0,1)^2$, with distinct $u$
coordinates and distinct $v$ coordinates. Let $m\ge16$, $r=1/m$,
and assume that no coordinate lies on a grid line and every open cell
\[
 (jr,(j+1)r)\times(kr,(k+1)r),\qquad 0\le j,k<m,
\]
contains a point. Put $J=[4r,1-4r]^2$ and
$B=\{i:x_i\notin J\}$. A realizer is a pair of strict total orders
$(L_1,L_2)$ whose intersection is the sampled chronology $P$.
For any strict order $L$, set
\[
 \rank_L(i)=1+\#\{j:L(j,i)\}.
\]
Let $\rank_u,\rank_v$ denote the coordinate ranks. Define empirical
marginals with inclusive thresholds, for example
$F_{v,n}(t)=n^{-1}\#\{i:v_i\le t\}$.

\begin{theorem}[Finite realizer rank rigidity]\label{thm:finite}
In addition to occupancy, suppose both empirical marginal errors on
$[0,1]$ are at most $2r$ and $\#B\le20rn$. For every realizer
$(L_1,L_2)$ there exists one exchange of these two orders, chosen for
the entire sample, such that for every $i$ with $x_i\in J$,
\[
 |\rank_{L_1}(i)-\rank_u(i)|\le30rn,\qquad
 |\rank_{L_2}(i)-\rank_v(i)|\le30rn.
\]
Here $L_1,L_2$ denote the orders after the chosen exchange.
\end{theorem}

\subsection{Orientation forcing}
For a realizer define
\[
 Q(x,y)\quad\Longleftrightarrow\quad L_1(x,y)\text{ and }L_2(y,x).
\]
Every incomparable pair has precisely one $Q$ orientation. The relation
$Q$ is transitive: two successive comparisons hold in $L_1$, and their
reverses hold in $L_2$. Thus their endpoints are incomparable and
have the corresponding $Q$ comparison.

If the incomparability edges $xy$ and $xz$ share $x$, while $y$ and $z$
are comparable in $P$, their orientations at $x$ agree. Otherwise a
two-edge $Q$ path would make $y,z$ incomparable by transitivity.
We use this elementary forcing rule for both the latent realizer and
any proposed realizer.

\subsection{A common anchor fixes separated pairs}
Occupancy supplies points in the following cells:
\begin{align*}
 A&\in(r,2r)\times(1-2r,1-r),\\
 B_0&\in(1-2r,1-r)\times(r,2r),\\
 W&\in(1-r,1)\times(1-3r,1-2r).
\end{align*}
The edges $AB_0$ and $AW$ are incomparability edges and $B_0\prec W$.
Take any interior incomparable $x,y$ with
$u_x<u_y$ and $v_x-v_y\ge3r$.
The interval $(v_y,v_x)$ contains an open grid interval: take
$k=\lfloor v_y/r\rfloor+1$; then $kr>v_y$ and
$(k+1)r\le v_y+2r<v_x$. Interior membership ensures an admissible
row. Occupancy in its first-column cell gives
\[
 z\in(0,r)\times(kr,(k+1)r).
\]
The coordinate inequalities imply
$z\prec x$, $z\prec A$, $y\prec W$, and $B_0\prec W$.
The sequence of incomparability edges
\[
 xy,\quad zy,\quad Ay,\quad AW,\quad AB_0
\]
has comparable opposite endpoints at each successive shared vertex.
Specifically the comparable pairs are $(z,x)$, $(z,A)$, $(y,W)$,
and $(B_0,W)$. The forcing rule therefore gives
\[
 Q(x,y)\ \Longleftrightarrow\ Q(z,y)
 \ \Longleftrightarrow\ Q(A,y)
 \ \Longleftrightarrow\ Q(A,W)
 \ \Longleftrightarrow\ Q(A,B_0).
\]
Exchanging $L_1,L_2$ reverses $Q$. Choose this exchange once so that
$Q(A,B_0)$ agrees with the latent $u$ order. All separated interior
incomparable pairs then agree. Comparable pairs agree automatically
in both realizers. No prime-graph or exact-uniqueness assumption is used.

\subsection{Counting possible disagreements}
For a fixed interior $x_i$, any disagreement in either order can only
involve a point outside $J$ or one with $|v_j-v_i|<3r$.
Set $a=v_i-3r$ and $b=v_i+3r$; both belong to $[0,1]$.
The strip is contained in $\{j:a<v_j\le b\}$, so the marginal bound gives
\[
 \#\{j:|v_j-v_i|<3r\}/n
 \le F_{v,n}(b)-F_{v,n}(a)\le(b-a)+4r=10r.
\]
At most $20rn+10rn$ partners can disagree. The absolute difference
of two predecessor counts is at most the number of points on which
their predecessor indicators differ. This proves both rank bounds,
with the same exchange, and establishes Theorem~\ref{thm:finite}.

\section{Order-only reconstruction and probability}\label{sec:probability}
For each labeled dimension-two poset $P$, fix one of the finitely many
pairs of total orders with intersection $P$, independently of the density.
For example, one may use the lexicographically first pair with labels
breaking ties. The Lean construction uses classical choice; the proof
applies to every realizer and does not depend on this choice. Define $A_n(P)$ as the
empirical cumulative distribution function of
\[
 \left(\frac{\rank_{L_1}(i)}n,\frac{\rank_{L_2}(i)}n\right),\qquad1\le i\le n.
\]
For other order codes define $A_n$ to be the zero function. This is a single function
of the observed order; grid points and anchors appear only in its
accuracy proof. No computational efficiency is asserted.

Write $C_\rho(s,t)=\int_{[0,s]\times[0,t]}\rho$ and let $C_n$ be the
latent empirical cumulative function. Suppose occupancy holds and
$|C_n-C_\rho|\le r$ at every grid vertex. By monotonicity between
vertices and the uniform population marginals:
\begin{enumerate}
\item each empirical marginal has error at most $2r$ at every threshold;
\item the full empirical cumulative error is at most $3r$;
\item the fraction outside $J$ is at most $20r$, since each of the
four marginal tails is bounded by $4r+r$.
\end{enumerate}
Coordinate ties and grid-line points have probability zero.
The finite theorem applies. Distinct coordinates and inclusive thresholds
give the exact identity $\rank_u(i)/n=F_{u,n}(u_i)$, and likewise for $v$;
there is no additional $1/n$ term. Each latent normalized rank differs
from its coordinate by at most $2r$, so every interior reconstructed point
is within $32r$ in each coordinate after one exchange.

Write $q_i$ for the aligned reconstructed point, $b=\#B/n\le20r$,
and $d=32r$. For $s,t\in[0,1]$ put
$s_- =\max(0,s-d)$, $s_+=\min(1,s+d)$ and similarly for $t$.
Interior-point indicator inclusions give, simultaneously for all thresholds,
\[
 C_n(s_-,t_-)-b\le A_n^S(s,t)\le C_n(s_+,t_+)+b.
\]
If a lower threshold clips to zero, its latent rectangle contains no
points because latent coordinates are strictly positive. Upper clipping
uses the unit-square bounds. This includes reconstructed ranks equal
to one and values $d>1$. Uniform population marginals bound the mass
change from moving the two thresholds by $2d=64r$; the latent empirical
replacement costs $3r$. Each direction uses one boundary penalty $b$.
Consequently
\begin{equation}\label{eq:reconstruction}
 \min\{\norm{A_n(P)-C_\rho},\norm{A_n(P)-C_{\rho^T}}\}\le87r.
\end{equation}
The transpose of a cumulative function exchanges its arguments and
preserves the supremum norm.

For $n\ge16^4$ choose $m=\lfloor n^{1/4}\rfloor$. Every cell has
mass at least $1/(2m^2)$, giving an occupancy failure bound
$m^2\exp(-n/(2m^2))$. At each grid vertex, the empirical mass is an
average of independent Bernoulli variables. Hoeffding's
inequality~\cite{Hoeffding} and a union bound give failure probability
$2(m+1)^2\exp(-2n/m^2)$ for vertex error at most $r$.
Independence across vertices or overlapping suborders is not assumed.
Since $n\ge m^4$ and $e^t\ge t^2/2$,
\begin{align*}
 m^2e^{-n/(2m^2)}+2(m+1)^2e^{-2n/m^2}
 &\le m^2e^{-m^2/2}+2(m+1)^2e^{-2m^2}\\
 &\le\frac8{m^2}+\frac{(m+1)^2}{m^4}
 \le\frac{12}{m^2}\le\frac{12}{256}<\frac1{10}.
\end{align*}
Also $m\ge n^{1/4}/2$. Thus~\eqref{eq:reconstruction} holds with
probability at least $9/10$, uniformly in $\rho\in\K$, with
\begin{equation}\label{eq:an}
 a_n=174n^{-1/4}
\end{equation}
in place of $87r$.

\section{From finite laws to coefficient error}\label{sec:inverse}
\begin{lemma}[Cumulative interpolation]\label{lem:interpolation}
For $\rho,\sigma\in\K$ and $\epsilon=\norm{C_\rho-C_\sigma}$,
\[
 \norm{\rho-\sigma}^3\le2048\epsilon.
\]
\end{lemma}
\begin{proof}
Put $f=\rho-\sigma$ and $\delta=\norm f\le1$. The function $f$ is
$4$-Lipschitz. If $\delta>0$, take a maximizer $x$ of $|f|$ and
$h=\delta/16$. In each coordinate choose a length-$h$ interval
extending from $x$ towards the half of $[0,1]$ with more room.
Their product $R$ lies in $D$. For $z\in R$,
$|z-x|_2\le2h$, so $|f(z)-f(x)|\le8h=\delta/2$.
The sign is fixed and
\[
 \left|\int_R f\right|\ge(\delta/2)h^2=\delta^3/512.
\]
Inclusion--exclusion expresses this integral by four corner values
of $C_\rho-C_\sigma$, of total absolute value at most $4\epsilon$.
The result follows; $\delta=0$ is immediate.
\end{proof}

Let
\[
 \epsilon_* =\min\{\norm{C_\rho-C_\sigma},
                         \norm{C_\rho-C_{\sigma^T}}\},\qquad
 t_n=\TV(\nu_n(\rho),\nu_n(\sigma)).
\]
Transposition is a supremum-norm isometry, so the minimum of the four
cross-distances between the two transpose orbits is exactly $\epsilon_*$.
The selector has a finite order domain; every output preimage used below
is therefore measurable without a function-space measurability assumption.
If $\epsilon_*>2a_n$, the sets of outputs within $a_n$ of the two
transpose orbits are disjoint. Each law places at least $9/10$ of
its mass on its own good set and at most $1/10$ on the other's.
Taking the preimage of a good set under the same order-only map
$A_n$ yields $t_n\ge4/5$. Therefore $t_n<4/5$ implies
$\epsilon_*\le2a_n$. Applying Lemma~\ref{lem:interpolation} to
the minimizing orientation yields
$\dc\le(4096a_n)^{1/3}$. If $t_n\ge4/5$, the elementary
$\dc\le1$ gives $\dc\le(5/4)t_n$. In both cases
\[
 \dc\le(4096a_n)^{1/3}+\tfrac54t_n.
\]
For $n\ge65536$, substitute~\eqref{eq:an} and use
$4096\cdot174=712704<100^3$. Setting $n=N$ and using
$t_N\le\Delta_N$ proves~\eqref{eq:main}. For $2\le N<65536$,
$\dc\le1<100N^{-1/12}$ proves it directly.

\begin{remark}[Magnitude and optimality]
The bare $N$ term becomes smaller than one only when $N>10^{24}$.
The bound establishes a polynomial modulus, not a useful finite sample
budget. No optimality of the $N^{-1/12}$ rate is asserted.
\end{remark}

\subsection{A logarithmic grid scale}
\begin{theorem}[Improved inverse estimate]\label{thm:improved}
For every $\rho,\sigma\in\K$ and integer $N\ge2$,
\[
 \dc(\rho,\sigma)\le130\bigl((\log N/N)^{1/6}+\Delta_N(\rho,\sigma)\bigr).
\]
The logarithm is natural. The observable, class and residual quotient
are those of Theorem~\ref{thm:main}.
\end{theorem}
\begin{proof}
For $n\ge65536$ set $x=\sqrt{n/(8\log n)}$, $m=\lfloor x\rfloor$,
and $r=1/m$. The function $t/\log t$ is increasing for $t>e$, since
its derivative is $(\log t-1)/(\log t)^2$. At the cutoff,
$\log65536=16\log2<16$, hence $x^2>512>16^2$.
Thus $m\ge16$ and $\log n\ge1$. Also $m\ge x-1\ge x/2$, so
\[
 m^2\le\frac{n}{8\log n},\qquad
 r\le2\sqrt{\frac{8\log n}{n}}.
\]
The deterministic reconstruction proof is unchanged. The two failure
terms in Section~\ref{sec:probability} satisfy
\begin{align*}
 m^2e^{-n/(2m^2)}&\le\frac{1}{8n^3\log n},\\
 2(m+1)^2e^{-2n/m^2}&\le\frac{n^{-15}}{\log n}.
\end{align*}
Indeed the exponents are at least $4\log n$ and $16\log n$;
$(m+1)^2\le4m^2$ supplies the second prefactor. Their sum is at most
$1/(8n^3)+n^{-15}<1/10$. Consequently success is at least $9/10$,
with cumulative error
\[
 a_n=174\sqrt8\sqrt{\log n/n}.
\]
The same orbit-separation and interpolation argument gives
\[
 \dc\le(712704\sqrt8)^{1/3}(\log n/n)^{1/6}+\tfrac54t_n.
\]
The coefficient is below 130: squaring the positive cubed comparison
reduces it to
\[
 8\cdot712704^2=4063575932928<4826809000000=130^6.
\]
Set $n=N$ and use $t_N\le\Delta_N$.

For $2\le N<65536$, $\log2\ge1/2$ follows by integrating $1/t$
on $[1,2]$. Thus $\log N/N\ge1/131072>(1/8)^6$, and the
displayed N term exceeds $130/8>1\ge\dc$. This proves the small-N
branch without a grid or concentration assertion there.
\end{proof}

The new estimate improves the asymptotic $N$ term but is not an optimal
rate or a practical sampling guarantee. It reuses the probability and
continuum bridges formalized for the original estimate. The logarithmic
grid choice and its resulting rate are not yet formalized.

\section{Observable labels and proper time}
\begin{proposition}[Labeled and unlabeled laws]
For each $k$, the total variation distance of the labeled order laws
equals that of their unlabeled order-isomorphism laws.
\end{proposition}
\begin{proof}
Exchangeability makes each labeled poset in a relabeling orbit $O$
have probability $p_O/|O|$. For the second law let the orbit mass be
$q_O$. Then labeled total variation is
\[
 \tfrac12\sum_O\sum_{P\in O}
       |p_O/|O|-q_O/|O||=\tfrac12\sum_O|p_O-q_O|,
\]
the unlabeled distance. Order duality is not included in relabeling.
\end{proof}
The label-based tie-breaking in $A_n$ therefore does not confer extra
information in these iid law comparisons.

For an absolutely continuous future-directed curve $\gamma=(u,v)$,
with nonnegative coordinate derivatives almost everywhere, define
\[
 L_\rho(\gamma)=\int\sqrt{2\rho(\gamma)u'v'}.
\]
Time separation is the supremum over such curves between the endpoints,
and is zero if no such future curve exists.
\begin{proposition}[Proper-time comparison]
There is one choice $S\in\{\mathrm{id},T\}$ such that, for all
endpoints $p,q\in D$,
\[
 |\tau_\rho(p,q)-\tau_\sigma(Sp,Sq)|\le\dc(\rho,\sigma).
\]
\end{proposition}
\begin{proof}
First set $\delta=\norm{\rho-\sigma}$. The density floor gives
\[
 |\sqrt{2\rho}-\sqrt{2\sigma}|
 =\frac{2|\rho-\sigma|}{\sqrt{2\rho}+\sqrt{2\sigma}}\le\delta.
\]
For a common monotone curve from $p$ to $q$, Cauchy--Schwarz gives
$\int\sqrt{u'v'}\le\sqrt{(q_u-p_u)(q_v-p_v)}\le1$.
Thus curve lengths differ by at most $\delta$. Both suprema are
finite, since every such length is at most $\sqrt3$, and the curve
classes agree. Taking suprema proves the identity-aligned bound.
Axis exchange preserves the product of derivatives and maps endpoints
and curves consistently. Choose the orientation attaining $\dc$.
\end{proof}

\section{Formal verification and reproducibility}\label{sec:formal}
The formalization uses the class in~\eqref{eq:class}: neighborhood
smoothness, bounds $1/2\le\rho\le3/2$, Euclidean Lipschitz constant two,
and both uniform marginals. The laws are pushforwards of the actual iid
density measure by the Boolean chronological relation. Observations
contain no coordinate ranks or latent positions. The unlabeled law is
the pushforward by directed-order isomorphism; order duality is not
part of that quotient.

The finite theorem appears in \path{QuantyraNullCone/Bridges.lean} as
\path{finite_realizer_rank_rigidity_specified}. Its existential Boolean
exchange precedes the universal event quantifier and is shared by both
rank bounds. The original inverse estimate appears as
\path{InDensityClass.full_inverse} in \path{InverseRate.lean}, and
\path{InDensityClass.full_inverse_unlabeled} in \path{Unlabeled.lean}.
These theorems assume only membership of both densities in $\K$ and
$N\ge2$. Concentration, reconstruction, interpolation and exchangeability
are proved rather than supplied as additional hypotheses.

The source modules, all under \path{QuantyraNullCone/}, cover the
following steps:
\begin{itemize}
\item \path{Model.lean}, \path{Measures.lean} and \path{DensityCDF.lean}
define the model and derive probability normalization, uniform marginals
and cumulative threshold stability.
\item \path{Cumulative.lean} and \path{GridAccuracy.lean} prove the
all-threshold $87r$ reconstruction estimate from the finite rank theorem.
\path{Sampling.lean}, \path{Coordinates.lean}, \path{Concentration.lean},
\path{GoodSamples.lean} and \path{ProbabilityRate.lean} discharge occupancy,
null-tie, concentration and integer-grid obligations, giving success
probability at least $9/10$ with error $174n^{-1/4}$ for $n\ge65536$.
\item \path{DensityInterpolation.lean} proves the coefficient bound
$\delta^3\le2048\epsilon$. \path{OrderSelector.lean}, \path{FiniteTV.lean}
and \path{Transpose.lean} prove the common order-only selection, total
variation separation and transpose identities. \path{InverseRate.lean}
handles the original all-$N$ estimate, including small $N$ and large TV.
\item \path{Exchangeability.lean}, \path{QuotientTV.lean} and
\path{Unlabeled.lean} derive iid relabeling invariance and exact equality
of labeled and unlabeled TV and $\Delta_N$. The all-size conclusion is
\path{InDensityClass.all_law_identifiability} in \path{Identifiability.lean}
and \path{InDensityClass.all_unlabeled_law_identifiability} in
\path{Unlabeled.lean}. Equality holds everywhere on $D$ after one global
identity or axis exchange; no equality outside $D$ is asserted.
\end{itemize}

Lean~4.30.0 and mathlib commit
\texttt{c5ea00351c28e24afc9f0f84379aa41082b1188f} are pinned.
The exact proof and evidence revision is
\href{https://github.com/Quantyra/quantyra-null-cone-lean/tree/b7d762acc9c10ca881f8366f545f3998b0528448}{\texttt{b7d762acc9c10ca881f8366f545f3998b0528448}}.
Authoritative verification ran on GCP under run identifier
\path{space-unlabeled-20261007T123238Z-759aaa}. The complete root build
finished successfully (2907 jobs), with 68 selected exact-type and
axiom audits and zero compiler warnings. Source hashes were checked
before and after execution, and all nine pinned dependency revisions
and their tracked trees were verified. The
\href{https://github.com/Quantyra/quantyra-null-cone-lean/tree/b7d762acc9c10ca881f8366f545f3998b0528448/evidence/gcp/space-unlabeled-20261007T123238Z-759aaa}{retained evidence directory}
contains immutable inputs, the manifest, raw logs and receipt.
The audited dependencies use only \texttt{propext},
\texttt{Classical.choice} and \texttt{Quot.sound}; no additional axioms
or admitted proofs are used. Hosted CI is supplementary to this GCP
verification. Reproduction instructions and the statement-to-source
mapping are in \path{INTEGRITY.md} and \path{notes/lean-verification.md}.

The logarithmic-grid estimate of Theorem~\ref{thm:improved} and the
proper-time comparison remain prose proofs. Their separate Lean
formalization is optional future work. Kernel checking does not establish
originality or replace scrutiny of the model and its interpretation.
Supplementary finite sanity checks cover 36 exact-coordinate grids
with 865445 separated edges and all 120 five-event coordinate permutations
with 771 valid realizers. The logarithmic-grid arithmetic check covers
612 sample sizes and 173 integer floor transitions. These finite checks
complement the proofs without certifying the optimized theorem in Lean.

\section{Discussion and research provenance}
The proof uses a fixed gauge and uniform positive density to separate
small geometric ambiguity from large-scale orientation. Extending
the class, finding sharper or optimal rates, controlling derivatives,
and formalizing the logarithmic-grid and proper-time consequences are
separate questions. This draft
does not claim to settle the general causal-set Hauptvermutung.

The work was developed with OpenAI Codex assistance under Daniel Eric
Fredriksen's research direction. AI-assisted derivation and writing are
disclosed. The open-source workflow supports decentralized informal
feedback, reuse and downstream testing; formal specialist review is
not a publication prerequisite. The repository retains the derivation, exact-source
comparisons and verification provenance. The candidate contribution
is the explicit restricted inverse estimate and the occupied-grid
quantification supporting its uniform use over nonconstant densities.
Earlier qualitative forcing and stronger flat-model rank recovery are
credited above; broader priority assessment remains provisional.

\begin{thebibliography}{9}
\bibitem{Janson} S. Janson, \emph{Poset limits and exchangeable random
posets}, arXiv:0902.0306. \url{https://arxiv.org/abs/0902.0306}.
\bibitem{Braun} M. Braun, \emph{Spacetime reconstruction by order and
number}, arXiv:2507.01907 (2025).
\url{https://arxiv.org/abs/2507.01907}.
\bibitem{Henson} J. Henson, \emph{Constructing an interval of Minkowski
space from a causal set}, arXiv:gr-qc/0601069 (2006).
\url{https://arxiv.org/abs/gr-qc/0601069}.
\bibitem{KlavikZeman} P. Klav\'ik and P. Zeman,
\emph{Automorphism Groups of Comparability Graphs},
arXiv:1506.05064. \url{https://arxiv.org/abs/1506.05064}.
\bibitem{Winkler1985} P. Winkler, \emph{Random orders},
Order \textbf{1}, 317--331 (1985).
\url{https://doi.org/10.1007/BF00582738}.
\bibitem{Winkler} P. Winkler, \emph{Random orders of dimension 2},
Order \textbf{7}, 329--339 (PDF dated 1991; publisher metadata 1990).
\url{https://doi.org/10.1007/BF00383197}.
\bibitem{Hoeffding} W. Hoeffding, \emph{Probability inequalities for
sums of bounded random variables}, Journal of the American Statistical
Association \textbf{58}(301), 13--30 (1963).
\url{https://doi.org/10.1080/01621459.1963.10500830}.
\end{thebibliography}
\end{document}
